
\documentclass{article}
\usepackage{colortbl}
\usepackage{makecell}
\usepackage{multirow}
\usepackage{supertabular}

\begin{document}

\newcounter{utterance}

\centering \large Interaction Transcript for game `dond', experiment `semi\_it', episode 12 with gpt{-}5{-}2025{-}08{-}07{-}no{-}reasoning{-}t1.0.
\vspace{24pt}

{ \footnotesize  \setcounter{utterance}{1}
\setlength{\tabcolsep}{0pt}
\begin{supertabular}{c@{$\;$}|p{.15\linewidth}@{}p{.15\linewidth}p{.15\linewidth}p{.15\linewidth}p{.15\linewidth}p{.15\linewidth}}
    \# & $\;$A & \multicolumn{4}{c}{Game Master} & $\;\:$B\\
    \hline

    \theutterance \stepcounter{utterance}  
    & & \multicolumn{4}{p{0.6\linewidth}}{
        \cellcolor[rgb]{0.9,0.9,0.9}{
            \makecell[{{p{\linewidth}}}]{
                \texttt{\tiny{[P1$\langle$GM]}}
                \texttt{State giocando a un gioco di negoziazione in cui dovete accordarvi su come dividere una serie di oggetti tra voi e un altro giocatore.} \\
\\ 
\texttt{Regole:} \\
\texttt{(a) A Lei e all'altro giocatore viene dato un insieme di oggetti. Ognuno di voi riceve anche una funzione di valore segreta, che rappresenta il valore di ciascun tipo di oggetto.} \\
\texttt{(b) Si scambiano messaggi con l'altro giocatore per concordare chi si aggiudica gli oggetti. Potete inviare un massimo di 5 messaggi ciascuno, oppure terminare in anticipo facendo una proposta segreta in qualsiasi momento.} \\
\texttt{(c) A ciascuno di voi viene chiesto di inviare una proposta segreta indicando gli oggetti che desiderate, formattata tra parentesi quadre come segue: "{[}Proposta: <numero> <nome oggetto>, <numero> <nome oggetto>, <...>{]}".} \\
\texttt{(d) Se le vostre proposte sono complementari, cioè ci sono abbastanza oggetti per soddisfare entrambe le proposte, a ciascun giocatore viene assegnato un punteggio basato sulla somma dei valori degli oggetti ricevuti. In caso contrario, entrambi ricevono zero punti.} \\
\texttt{(e) L'obiettivo è massimizzare il proprio punteggio.} \\
\\ 
\texttt{Cominciamo.} \\
\\ 
\texttt{L'insieme degli oggetti disponibili è:} \\
\\ 
\texttt{1 lampada, 3 ditali, 1 corda} \\
\\ 
\texttt{La funzione valore segreta è:} \\
\\ 
\texttt{lampada: 3, ditale: 2, corda: 1} \\
\\ 
\texttt{IMPORTANTE: i vostri messaggi, a meno che non si tratti di una proposta segreta, vengono trasmessi direttamente all'altro giocatore, quindi non includete alcuna risposta alle regole o testo di annuncio del vostro messaggio. Per fare una proposta segreta, utilizzate il formato indicato. Non utilizzare delle parentesi quadre quando si comunica all'altro giocatore, altrimenti verrà interpretata come una proposta segreta.} \\
            }
        }
    }
    & \\ \\

    \theutterance \stepcounter{utterance}  
    & \multicolumn{4}{p{0.6\linewidth}}{
        \cellcolor[rgb]{0.9,0.9,0.9}{
            \makecell[{{p{\linewidth}}}]{
                \texttt{\tiny{[P1$\rangle$GM]}}
                \texttt{Ciao! Ecco cosa propongo per iniziare: per me i ditali valgono abbastanza, la lampada ancora di più e la corda un po’ meno. Ti va se mi dici quali oggetti per te sono più importanti? Possiamo cercare uno scambio dove ognuno prende ciò che valuta di più.} \\
            }
        }
    }
    & & \\ \\

    \theutterance \stepcounter{utterance}  
    & & & \multicolumn{2}{p{0.3\linewidth}}{
        \cellcolor[rgb]{0.9,0.9,0.9}{
            \makecell[{{p{\linewidth}}}]{
                \texttt{\tiny{[GM$|$GM]}}
                \texttt{continue} \\
            }
        }
    }
    & & \\ \\

    \theutterance \stepcounter{utterance}  
    & & & \multicolumn{2}{p{0.3\linewidth}}{
        \cellcolor[rgb]{0.9,0.9,0.9}{
            \makecell[{{p{\linewidth}}}]{
                \texttt{\tiny{[GM$|$GM]}}
                \texttt{Ciao! Ecco cosa propongo per iniziare: per me i ditali valgono abbastanza, la lampada ancora di più e la corda un po’ meno. Ti va se mi dici quali oggetti per te sono più importanti? Possiamo cercare uno scambio dove ognuno prende ciò che valuta di più.} \\
            }
        }
    }
    & & \\ \\

    \theutterance \stepcounter{utterance}  
    & & \multicolumn{4}{p{0.6\linewidth}}{
        \cellcolor[rgb]{0.9,0.9,0.9}{
            \makecell[{{p{\linewidth}}}]{
                \texttt{\tiny{[GM$\rangle$P2]}}
                \texttt{State giocando a un gioco di negoziazione in cui dovete accordarvi su come dividere una serie di oggetti tra voi e un altro giocatore.} \\
\\ 
\texttt{Regole:} \\
\texttt{(a) A Lei e all'altro giocatore viene dato un insieme di oggetti. Ognuno di voi riceve anche una funzione di valore segreta, che rappresenta il valore di ciascun tipo di oggetto.} \\
\texttt{(b) Si scambiano messaggi con l'altro giocatore per concordare chi si aggiudica gli oggetti. Potete inviare un massimo di 5 messaggi ciascuno, oppure terminare in anticipo facendo una proposta segreta in qualsiasi momento.} \\
\texttt{(c) A ciascuno di voi viene chiesto di inviare una proposta segreta indicando gli oggetti che desiderate, formattata tra parentesi quadre come segue: "{[}Proposta: <numero> <nome oggetto>, <numero> <nome oggetto>, <...>{]}".} \\
\texttt{(d) Se le vostre proposte sono complementari, cioè ci sono abbastanza oggetti per soddisfare entrambe le proposte, a ciascun giocatore viene assegnato un punteggio basato sulla somma dei valori degli oggetti ricevuti. In caso contrario, entrambi ricevono zero punti.} \\
\texttt{(e) L'obiettivo è massimizzare il proprio punteggio.} \\
\\ 
\texttt{Cominciamo.} \\
\\ 
\texttt{L'insieme degli oggetti disponibili è:} \\
\\ 
\texttt{1 lampada, 3 ditali, 1 corda} \\
\\ 
\texttt{La funzione valore segreta è:} \\
\\ 
\texttt{lampada: 3, ditale: 1, corda: 4} \\
\\ 
\texttt{IMPORTANTE: i vostri messaggi, a meno che non si tratti di una proposta segreta, vengono trasmessi direttamente all'altro giocatore, quindi non includete alcuna risposta alle regole o testo di annuncio del vostro messaggio. Per fare una proposta segreta, utilizzate il formato indicato. Non utilizzare delle parentesi quadre quando si comunica all'altro giocatore, altrimenti verrà interpretata come una proposta segreta.} \\
\\ 
\\ 
\texttt{Ciao! Ecco cosa propongo per iniziare: per me i ditali valgono abbastanza, la lampada ancora di più e la corda un po’ meno. Ti va se mi dici quali oggetti per te sono più importanti? Possiamo cercare uno scambio dove ognuno prende ciò che valuta di più.} \\
            }
        }
    }
    & \\ \\

    \theutterance \stepcounter{utterance}  
    & & & \multicolumn{4}{p{0.6\linewidth}}{
        \cellcolor[rgb]{0.9,0.9,0.9}{
            \makecell[{{p{\linewidth}}}]{
                \texttt{\tiny{[GM$\langle$P2]}}
                \texttt{Grazie! Per me la lampada è abbastanza importante, i ditali hanno un valore moderato e la corda è quella che valuto di più. Ti andrebbe, se per te i ditali contano, di prendere tu i ditali, mentre io prendo la corda? Possiamo poi discutere della lampada: tu quanto la valuti rispetto ai ditali?} \\
            }
        }
    }
     \\ \\

    \theutterance \stepcounter{utterance}  
    & & & \multicolumn{2}{p{0.3\linewidth}}{
        \cellcolor[rgb]{0.9,0.9,0.9}{
            \makecell[{{p{\linewidth}}}]{
                \texttt{\tiny{[GM$|$GM]}}
                \texttt{continue} \\
            }
        }
    }
    & & \\ \\

    \theutterance \stepcounter{utterance}  
    & & & \multicolumn{2}{p{0.3\linewidth}}{
        \cellcolor[rgb]{0.9,0.9,0.9}{
            \makecell[{{p{\linewidth}}}]{
                \texttt{\tiny{[GM$|$GM]}}
                \texttt{Grazie! Per me la lampada è abbastanza importante, i ditali hanno un valore moderato e la corda è quella che valuto di più. Ti andrebbe, se per te i ditali contano, di prendere tu i ditali, mentre io prendo la corda? Possiamo poi discutere della lampada: tu quanto la valuti rispetto ai ditali?} \\
            }
        }
    }
    & & \\ \\

    \theutterance \stepcounter{utterance}  
    & & \multicolumn{4}{p{0.6\linewidth}}{
        \cellcolor[rgb]{0.9,0.9,0.9}{
            \makecell[{{p{\linewidth}}}]{
                \texttt{\tiny{[P1$\langle$GM]}}
                \texttt{Grazie! Per me la lampada è abbastanza importante, i ditali hanno un valore moderato e la corda è quella che valuto di più. Ti andrebbe, se per te i ditali contano, di prendere tu i ditali, mentre io prendo la corda? Possiamo poi discutere della lampada: tu quanto la valuti rispetto ai ditali?} \\
            }
        }
    }
    & \\ \\

    \theutterance \stepcounter{utterance}  
    & \multicolumn{4}{p{0.6\linewidth}}{
        \cellcolor[rgb]{0.9,0.9,0.9}{
            \makecell[{{p{\linewidth}}}]{
                \texttt{\tiny{[P1$\rangle$GM]}}
                \texttt{Grazie, chiarissimo. Per me la lampada vale più dei ditali, e i ditali più della corda. Visto che tu tieni molto alla corda, per me va bene che tu prenda la corda.} \\
\\ 
\texttt{Sulla lampada: per me vale più di un ditale singolo, ma non più di due ditali messi insieme. Quindi potremmo fare così:} \\
\texttt{{-} Tu prendi la corda} \\
\texttt{{-} Io prendo la lampada} \\
\texttt{{-} Dividiamo i 3 ditali: io ne prenderei 2 e tu 1} \\
\\ 
\texttt{Se ti va, possiamo anche invertire la divisione dei ditali (tu 2 e io 1) se in cambio la lampada va a me. Fammi sapere quale delle due varianti preferisci, oppure proponi tu una ripartizione.} \\
            }
        }
    }
    & & \\ \\

    \theutterance \stepcounter{utterance}  
    & & & \multicolumn{2}{p{0.3\linewidth}}{
        \cellcolor[rgb]{0.9,0.9,0.9}{
            \makecell[{{p{\linewidth}}}]{
                \texttt{\tiny{[GM$|$GM]}}
                \texttt{continue} \\
            }
        }
    }
    & & \\ \\

    \theutterance \stepcounter{utterance}  
    & & & \multicolumn{2}{p{0.3\linewidth}}{
        \cellcolor[rgb]{0.9,0.9,0.9}{
            \makecell[{{p{\linewidth}}}]{
                \texttt{\tiny{[GM$|$GM]}}
                \texttt{Grazie, chiarissimo. Per me la lampada vale più dei ditali, e i ditali più della corda. Visto che tu tieni molto alla corda, per me va bene che tu prenda la corda.} \\
\\ 
\texttt{Sulla lampada: per me vale più di un ditale singolo, ma non più di due ditali messi insieme. Quindi potremmo fare così:} \\
\texttt{{-} Tu prendi la corda} \\
\texttt{{-} Io prendo la lampada} \\
\texttt{{-} Dividiamo i 3 ditali: io ne prenderei 2 e tu 1} \\
\\ 
\texttt{Se ti va, possiamo anche invertire la divisione dei ditali (tu 2 e io 1) se in cambio la lampada va a me. Fammi sapere quale delle due varianti preferisci, oppure proponi tu una ripartizione.} \\
            }
        }
    }
    & & \\ \\

    \theutterance \stepcounter{utterance}  
    & & \multicolumn{4}{p{0.6\linewidth}}{
        \cellcolor[rgb]{0.9,0.9,0.9}{
            \makecell[{{p{\linewidth}}}]{
                \texttt{\tiny{[GM$\rangle$P2]}}
                \texttt{Grazie, chiarissimo. Per me la lampada vale più dei ditali, e i ditali più della corda. Visto che tu tieni molto alla corda, per me va bene che tu prenda la corda.} \\
\\ 
\texttt{Sulla lampada: per me vale più di un ditale singolo, ma non più di due ditali messi insieme. Quindi potremmo fare così:} \\
\texttt{{-} Tu prendi la corda} \\
\texttt{{-} Io prendo la lampada} \\
\texttt{{-} Dividiamo i 3 ditali: io ne prenderei 2 e tu 1} \\
\\ 
\texttt{Se ti va, possiamo anche invertire la divisione dei ditali (tu 2 e io 1) se in cambio la lampada va a me. Fammi sapere quale delle due varianti preferisci, oppure proponi tu una ripartizione.} \\
            }
        }
    }
    & \\ \\

    \theutterance \stepcounter{utterance}  
    & & & \multicolumn{4}{p{0.6\linewidth}}{
        \cellcolor[rgb]{0.9,0.9,0.9}{
            \makecell[{{p{\linewidth}}}]{
                \texttt{\tiny{[GM$\langle$P2]}}
                \texttt{Grazie, mi sembra una base equa. Dato che per me la corda è la priorità, sono d’accordo a prenderla. Tra le due varianti che proponi, preferisco questa:} \\
\texttt{{-} Tu prendi la lampada} \\
\texttt{{-} Io prendo la corda} \\
\texttt{{-} Ditali: tu 1, io 2} \\
\\ 
\texttt{Se per te va bene così, confermiamo; altrimenti posso accettare anche l’altra divisione dei ditali pur di mantenere la lampada a te e la corda a me.} \\
            }
        }
    }
     \\ \\

    \theutterance \stepcounter{utterance}  
    & & & \multicolumn{2}{p{0.3\linewidth}}{
        \cellcolor[rgb]{0.9,0.9,0.9}{
            \makecell[{{p{\linewidth}}}]{
                \texttt{\tiny{[GM$|$GM]}}
                \texttt{continue} \\
            }
        }
    }
    & & \\ \\

    \theutterance \stepcounter{utterance}  
    & & & \multicolumn{2}{p{0.3\linewidth}}{
        \cellcolor[rgb]{0.9,0.9,0.9}{
            \makecell[{{p{\linewidth}}}]{
                \texttt{\tiny{[GM$|$GM]}}
                \texttt{Grazie, mi sembra una base equa. Dato che per me la corda è la priorità, sono d’accordo a prenderla. Tra le due varianti che proponi, preferisco questa:} \\
\texttt{{-} Tu prendi la lampada} \\
\texttt{{-} Io prendo la corda} \\
\texttt{{-} Ditali: tu 1, io 2} \\
\\ 
\texttt{Se per te va bene così, confermiamo; altrimenti posso accettare anche l’altra divisione dei ditali pur di mantenere la lampada a te e la corda a me.} \\
            }
        }
    }
    & & \\ \\

    \theutterance \stepcounter{utterance}  
    & & \multicolumn{4}{p{0.6\linewidth}}{
        \cellcolor[rgb]{0.9,0.9,0.9}{
            \makecell[{{p{\linewidth}}}]{
                \texttt{\tiny{[P1$\langle$GM]}}
                \texttt{Grazie, mi sembra una base equa. Dato che per me la corda è la priorità, sono d’accordo a prenderla. Tra le due varianti che proponi, preferisco questa:} \\
\texttt{{-} Tu prendi la lampada} \\
\texttt{{-} Io prendo la corda} \\
\texttt{{-} Ditali: tu 1, io 2} \\
\\ 
\texttt{Se per te va bene così, confermiamo; altrimenti posso accettare anche l’altra divisione dei ditali pur di mantenere la lampada a te e la corda a me.} \\
            }
        }
    }
    & \\ \\

    \theutterance \stepcounter{utterance}  
    & \multicolumn{4}{p{0.6\linewidth}}{
        \cellcolor[rgb]{0.9,0.9,0.9}{
            \makecell[{{p{\linewidth}}}]{
                \texttt{\tiny{[P1$\rangle$GM]}}
                \texttt{Confermo: va benissimo per me.} \\
\\ 
\texttt{{[}Proposta: 1 lampada, 1 ditale{]}} \\
\\ 
\texttt{(Considerando la tua, dovrebbe essere: 1 corda, 2 ditali per te. Se concordiamo, le proposte sono complementari.)} \\
            }
        }
    }
    & & \\ \\

    \theutterance \stepcounter{utterance}  
    & & & \multicolumn{2}{p{0.3\linewidth}}{
        \cellcolor[rgb]{0.9,0.9,0.9}{
            \makecell[{{p{\linewidth}}}]{
                \texttt{\tiny{[GM$|$GM]}}
                \texttt{proposal} \\
            }
        }
    }
    & & \\ \\

    \theutterance \stepcounter{utterance}  
    & & & \multicolumn{2}{p{0.3\linewidth}}{
        \cellcolor[rgb]{0.9,0.9,0.9}{
            \makecell[{{p{\linewidth}}}]{
                \texttt{\tiny{[GM$|$GM]}}
                [1, 1, 0]
            }
        }
    }
    & & \\ \\

    \theutterance \stepcounter{utterance}  
    & & \multicolumn{4}{p{0.6\linewidth}}{
        \cellcolor[rgb]{0.9,0.9,0.9}{
            \makecell[{{p{\linewidth}}}]{
                \texttt{\tiny{[GM$\rangle$P2]}}
                \texttt{L'altro giocatore ha presentato una proposta segreta. Ora tocca a Lei presentare una sua proposta.} \\
            }
        }
    }
    & \\ \\

    \theutterance \stepcounter{utterance}  
    & & & \multicolumn{4}{p{0.6\linewidth}}{
        \cellcolor[rgb]{0.9,0.9,0.9}{
            \makecell[{{p{\linewidth}}}]{
                \texttt{\tiny{[GM$\langle$P2]}}
                \texttt{{[}Proposta: 1 lampada, 1 ditale, 1 corda{]}} \\
            }
        }
    }
     \\ \\

    \theutterance \stepcounter{utterance}  
    & & & \multicolumn{2}{p{0.3\linewidth}}{
        \cellcolor[rgb]{0.9,0.9,0.9}{
            \makecell[{{p{\linewidth}}}]{
                \texttt{\tiny{[GM$|$GM]}}
                \texttt{proposal} \\
            }
        }
    }
    & & \\ \\

    \theutterance \stepcounter{utterance}  
    & & & \multicolumn{2}{p{0.3\linewidth}}{
        \cellcolor[rgb]{0.9,0.9,0.9}{
            \makecell[{{p{\linewidth}}}]{
                \texttt{\tiny{[GM$|$GM]}}
                [1, 1, 1]
            }
        }
    }
    & & \\ \\

    \theutterance \stepcounter{utterance}  
    & & & \multicolumn{2}{p{0.3\linewidth}}{
        \cellcolor[rgb]{0.9,0.9,0.9}{
            \makecell[{{p{\linewidth}}}]{
                \texttt{\tiny{[GM$|$GM]}}
                [[1, 1, 0], [1, 1, 1]]
            }
        }
    }
    & & \\ \\

\end{supertabular}
}

\end{document}
